\begin{center}\begin{tikzpicture}[font=\small,>=stealth,line join=round]
\draw[gray!40,rounded corners=3pt] (0,0) rectangle (14.2,8.5);\node at (7.1,7.9) {非零截距会怎样破坏原点对称？};
\begin{scope}[shift={(4.0,3.8)},x=1.2cm,y=0.9cm]\draw[->,gray] (-1.9,0)--(1.9,0) node[right]{$x$};\draw[->,gray] (0,-2.8)--(0,2.8) node[above]{$y$};\draw[blue,thick,domain=-1.05:1.05,samples=150] plot(\x,{(\x)^3+1});\draw[violet,dashed,domain=-1.05:1.05,samples=100] plot(\x,{(\x)^3-1});\fill[blue] (0,1) circle(2.2pt) node[above left=5pt] {$(0,1)$};\draw[violet,fill=white] (0,-1) circle(2.5pt);\node[below right=8pt,violet] at (0,-1){$(0,-1)$};\end{scope}\node[align=left,anchor=west] at (8.0,4.1){蓝实线：$y=x^3+1$\\[11pt]紫虚线：旋转半圈后的图像\\[11pt]两图不重合，不能是奇函数\\[11pt]只看 $f(0)=1\ne0$ 就能排除};
\node[text width=12.8cm,align=center] at (7.1,.65) {虚线为蓝图的原点对称像 $y=x^3-1$，不是原函数的另一部分。};
\end{tikzpicture}\end{center}
